\documentclass[%
ps,
%pdf,
%nocolorBG,
%colorBG,
%slideColor,
%slideBW,
%draft,
%frames
%azure
%contemporain
%nuancegris
%troispoints
%lignesbleues
%darkblue
%alienglow
%autumn
%blueline
]{prosper}
\usepackage{amsmath,alltt}

\usepackage{/usr/lib/R/share/texmf/Sweave}

\newcommand{\Ran}{\,\mathrm{Ran}\,}

%%% Things which go on all slides:
%% the caption. The default is the \title argument.
%% \slideCaption{Copula Package}
%% A logo. x and y coordinates measured in inches with
%% zero point at the lower left corner.
\Logo{\includegraphics{uilogo}}

\title{R Package Copula}
\author{Jun Yan}
\institution{Department of Statistics and Actuarial Science,\\
University of Iowa}
\slideCaption{Copula Package}


\begin{document}


\maketitle

\overlays{6}{
\begin{slide}{Backgrounds}
\begin{itemstep}
\item Copulas have become a widely used tool for 
modeling multivariate dependence in a variety of fields.

\item Software implementation is important in promoting
the application of copulas.

\item \texttt{Splus} has a collection of functions for copulas
modeling in the \texttt{finmetrics} module.

\begin{itemize}
\item Commercial
\item Bivariate copulas only
\end{itemize}

\item Need a platform for the development of 
copula methods and applications.

%\item \texttt{R} package \texttt{copula}.
\end{itemstep}
\end{slide}
}


\begin{slide}{Copulas}
\begin{itemize}
\item A copula $C$ is a multivariate distribution
whose margins are all uniform over $(0,1)$:
$$C(u_1,\ldots, u_p) = \Pr(U_1 \le u1, \ldots, U_p \le u_p).$$

\item Sklar's cannonical representaion theorem (1959): 
A multivariate joint distribution can be represented 
by its marginal distribution and a copula:
$$F(x_1, \ldots, x_p) = C[F_1(x_1), \ldots, F_p(x_p)].$$
The copula is unique if the margins are continuous.
Otherwise, only the subcopulas is uniquely determined
on $\Ran F_1 \times \cdots \times \Ran F_p$.

\item Sklar's Theorem (1959)

\end{itemize}
  
\end{slide}

\overlays{7}{
\begin{slide}{Why R}
\begin{itemstep}
\item Quote from \texttt{http://www.r-project.org}:
  ``\texttt{R} is a free software environment for 
  statistical computing and graphics.'' 

\item Open source.

\item Compiles and runs on a wide variety of UNIX platforms, 
Windows and MacOS. 

\item Cutting-edge development;
hundreds of contributed packages.

\item Excellent graphics.

\item Easy interface with lower level compiled code (C/C++, Fortran)

\item Active developer-user interaction.
\end{itemstep}
  
\end{slide}
}

\overlays{10}{
\begin{slide}{Features of the Copula Package}
\begin{itemstep}
\item Classes (S4) of commonly used copula families
\begin{itemize}
\item Elliptical copulas: normal, t, Clayton, Frank, and Gumbel
\item Archimedian copulas
\item Extreme value copulas (to be implemented)
\end{itemize}

\item Dimension can be greater than 2.
\item Methods
\begin{itemize}
\item density
\item distribution
\item random number generator
\end{itemize}
\item Graphics: perspective plot, contour plot.
\end{itemstep}

\end{slide}
}

\begin{slide}{Load the Package}
The package \texttt{copula} depends on contributed
packages \texttt{mvtnorm}, \texttt{scatterplot3d},
and package \texttt{sn}.

\begin{Schunk}
\begin{Sinput}
> library(copula)
\end{Sinput}
\end{Schunk}
  
\end{slide}

\begin{slide}{Copula Objects}
  
\end{slide}

\begin{slide}{Density Contour}
  
\end{slide}

\begin{slide}{Density Perspective}
  
\end{slide}


\end{document}
